% !TEX program = pdflatex
\documentclass[11pt,a4paper]{article}

\usepackage[a4paper,top=25mm,bottom=25mm,left=30mm,right=30mm]{geometry}
\usepackage[utf8]{inputenc}
\usepackage[T1]{fontenc}
\usepackage{lmodern}
\usepackage{microtype}
\microtypesetup{patch=none}
\usepackage{amsmath,amssymb,amsthm}
\usepackage{mathtools}
\usepackage{graphicx}
\usepackage{xcolor}
\usepackage{tikz}
\usetikzlibrary{arrows.meta,positioning}
\usepackage{booktabs}
\usepackage{multirow}
\usepackage{enumitem}
\usepackage{setspace}
\usepackage{caption}
\usepackage[numbers,sort&compress]{natbib}
\usepackage[colorlinks=true,linkcolor=blue!60!black,citecolor=green!40!black]{hyperref}
\usepackage[ruled,vlined]{algorithm2e}

\newtheorem{theorem}{Theorem}[section]
\newtheorem{lemma}[theorem]{Lemma}
\newtheorem{corollary}[theorem]{Corollary}
\newtheorem{definition}{Definition}[section]

\newcommand{\MI}{\mathrm{MI}}
\newcommand{\MIp}{\mathrm{MIp}}

\pagestyle{plain}

\begin{document}

\title{Machine-Provable Karnaugh Map Encoding of Protein Sequences:\\
	Boolean Circuit Extraction and Cross-Protein Coevolutionary\\
	Contact Prediction}

\author{Shuvam Banerji Seal\\
	\small Indian Institute of Science Education and Research Kolkata, Mohanpur 741246, India}

\date{}

\maketitle

\begin{abstract}
	The representation of protein sequences as Karnaugh maps through Gray-code
	encoding establishes a formal bridge between switching theory and the
	statistical mechanics of coevolution.  We present a computational framework in
	which the entire analytical chain---from amino-acid encoding through
	Karnaugh-map construction, Quine--McCluskey minimisation, and prime-implicant
	extraction to cross-protein contact prediction---is formalised in the Lean~4
	theorem prover and validated against native experimental structures from the
	PSICOV150 benchmark.  Our universal contact circuits, trained on consensus
	residue identities from approximately one hundred proteins, transfer to
	held-out proteins at a mean top-$L/5$ precision of 0.179 ($5.1\times$ random
	enrichment).  Fusing circuit-derived scores with continuous DCA rankings via
	rank aggregation yields a hybrid predictor that exceeds both individual
	methods ($p < 10^{-4}$), demonstrating for the first time that
	Boolean-minimisation features carry structural information orthogonal to
	standard covariation analysis.  Flipped (forbidden-combination) circuits
	generalise across proteins with 98.3\% specificity.  Gray-code adjacency is
	enriched among native contacts ($1.178\times$, $p < 10^{-77}$), and conserved
	column pairs are contact-enriched by $3.2\times$.  We further prove an
	information ceiling: any predictor that is a function of a residue-level
	feature map is bounded by that map's cell purity, which on consensus residue
	identity and separation is prec@$L/5 = 0.193$---so the trained circuit already
	attains $93\%$ of what its inputs permit, and the gap to direct-coupling
	analysis is informational rather than algorithmic.  Redirecting the same
	minimisation from sequence to \emph{conformational state} escapes that bound:
	across 480 structure-based-model trajectories the extracted rules link
	contacts that are $12.5\times$ enriched for contact-map adjacency
	($p < 10^{-79}$), recovering local folding cooperativity, and under dense
	re-sampling they predict a contact's state a thousand integration steps ahead
	in every protein tested.  All software, formal
	proofs, and result files are publicly available.
\end{abstract}

% ── Main text ──
% ═══════════════════════════════════════════════════════════
% §1 Introduction
% ═══════════════════════════════════════════════════════════
\section{Introduction}
\label{sec:intro}

Residue--residue contact prediction from multiple sequence alignments has
transformed protein structure determination~\cite{weigt2009,morcos2011,
	jones2012,kamisetty2013}.  The dominant paradigm, Direct Coupling Analysis
(DCA), models the joint distribution of amino-acid sequences through a Potts
Hamiltonian whose coupling parameters are estimated by inverting a covariance
matrix computed from deep alignments~\cite{morcos2011}.  Sparse inverse
covariance estimation, as implemented in PSICOV~\cite{jones2012}, achieves a
mean top-$L/5$ precision of approximately 0.73 on standard benchmarks.
Pseudo-likelihood maximisation~\cite{ekeberg2013} and related methods further
improve upon mean-field approximations.  More recently, end-to-end deep
learning approaches such as AlphaFold2~\cite{jumper2021} have largely
superseded standalone contact prediction by integrating coevolutionary signal
with structural priors inside attention-based neural architectures.

Despite this success, the encoding layer through which biological sequences
enter the statistical machinery remains opaque.  No prior work has
established machine-checkable proofs that the transformation from amino-acid
alphabet to the numerical representation used for inference preserves the
structural properties that downstream analysis requires.  This gap is not
merely academic: subtle encoding errors can silently corrupt results while
producing plausible-looking output.

Here we address this gap by constructing a formally verified bridge between
digital-logic design theory and protein coevolution analysis.  Our central
observation is that the Karnaugh map~\cite{karnaugh1953}, when combined with
reflected binary Gray-code ordering of amino acids~\cite{gray1953}, provides
a natural embedding of aligned sequences into a grid structure where adjacent
cells correspond to physicochemically similar residue pairs.  The
Quine--McCluskey algorithm~\cite{quine1952,mccluskey1956} then extracts
minimal sum-of-products expressions from thresholded frequency maps,
producing prime implicants that we interpret as coevolutionary constraints.

We formalise 118 theorems in Lean~4 covering Gray-code bijectivity, $k$-mer
indexing injectivity, amino-acid encoding properties, contact-map topology,
implicant-cover soundness, and padding safety.  We apply the verified pipeline
to 150 protein domains from the PSICOV benchmark comprising over 1.6 million
labelled observations against native experimental structures.  Our results
demonstrate that Boolean circuits transfer across proteins, that their
cell-rate prior improves DCA ranking when fused via rank aggregation, and
that negative-selection (forbidden-combination) circuits generalise with high
specificity.  We also report negative results honestly: an attention-style
tension mechanism does not improve prediction, iterative direct information
is unreliable on conserved-rich families, and current circuit precision falls
far short of the level required for template-free structure reconstruction.

All software, formal proofs, result files, and documentation are publicly
available (\texttt{github.com/E-Motioner-X-SBS/co-evolution-analysis}).

% ═══ Related Work ═══
\section{Related Work}
\label{sec:related}

\subsection{Coevolutionary contact prediction}
The inference of residue--residue contacts from multiple sequence alignments
(MSAs) has evolved through several generations of methodology.  Early
approaches based on mutual information~\cite{dunn2008} suffered from
phylogenetic and entropic biases that produced high false-positive rates.
Direct Coupling Analysis (DCA) addressed this by modelling the joint sequence
distribution through a Potts Hamiltonian whose coupling parameters capture
direct (rather than transitive) correlations~\cite{weigt2009}.  The mean-field
approximation of Morcos et al.\cite{morcos2011} made DCA computationally
tractable for large alignments by replacing the full partition function with a
covariance-based estimate.  Sparse inverse covariance estimation, as
implemented in PSICOV~\cite{jones2012}, further improved precision by
promoting sparsity in the precision matrix through $\ell_1$ regularisation.
Pseudo-likelihood maximisation~\cite{ekeberg2013} and fast belief propagation
implementations such as CCMpred~\cite{seemayer2014} subsequently refined the
inference procedure.  GREMLIN~\cite{kamisetty2013} combined DCA with
protein-structure priors.  EVmutation~\cite{hopf2017} extended the framework
to predict mutation effects on fitness rather than contacts alone.

More recently, end-to-end deep learning approaches such as
AlphaFold2~\cite{jumper2021} have integrated coevolutionary signal with
structural priors inside attention-based architectures, achieving near-experimental
accuracy for many targets.  However, these models produce dense internal
representations that resist decomposition into interpretable components.

\subsection{Formal methods in computational biology}
The application of interactive theorem provers to computational biology
remains limited.  Existing formalisations focus on algorithm correctness for
phylogenetic inference~\cite{lean4}, genome assembly graph properties, and
sequence alignment dynamic programming recurrences.  No prior work has
established machine-checked proofs that the encoding layer connecting amino-acid
alphabets to numerical representations preserves the structural properties
required for coevolutionary analysis.

\subsection{Karnaugh maps beyond digital logic}
Karnaugh maps were introduced as a graphical method for minimising Boolean
functions~\cite{karnaugh1953}.  Their application outside digital circuit
design has been explored in synthetic biology, where Boolean gates model gene
regulatory networks~\cite{marchisio2011}, and in the analysis of genetic code
symmetries~\cite{petoukhov2025}.  Our work extends this tradition by using
Karnaugh maps not merely as a visual representation but as an analytical
engine: the Quine--McCluskey minimiser extracts testable coevolutionary rules
from thresholded frequency data, and the resulting circuits are validated
against native protein structures.

\section{Gray-Code Encoding of Biological Alphabets}
\label{sec:theory-gray}

A reflected binary Gray code assigns to each non-negative integer $n$ the
value $g(n) = n \oplus (n \gg 1)$~\cite{gray1953}.  Consecutive integers map
to codewords differing in exactly one bit.  For nucleotides we adopt the
standard two-bit encoding A=0, C=1, G=3, T/U=2; transitions correspond to
single bit-flips under Gray transformation.

For amino acids, each residue receives its He~2012 integer
code~\cite{he2012}: A=0 through G=19, ordered into seven physicochemical
groups (aliphatic AILVM, aromatic FYW, sulfur MC, structure PG, polar STNQ,
negative DE, positive HKR).  The five-bit Gray image $g(\mathrm{code}(a))$
places each residue on the hypercube $Q_5$.  Among the $\binom{20}{2}=190$
unordered pairs, exactly 80 ordered pairs lie at Gray distance~1 (verified in
Lean~4 as \texttt{cc\_he\_dist1\_ordered\_eq\_80}).

Alignment gaps are mapped to state~20 and excluded from all statistical
calculations but retained in position arrays so that column identity is
preserved across all sequences.

\section{Karnaugh Maps and Quine--McCluskey Minimisation}
\label{sec:theory-kmap-qm}

For a pair of alignment columns $(i,j)$, a ternary Karnaugh map $K_{ij}$
assigns each cell of a $32\times32$ grid one of three values based on joint
residue frequencies: ON for observed mutations, OFF for never-observed pairs,
and don't-care for reference pairs and padding rows/columns 20--31.

The Quine--McCluskey algorithm~\cite{quine1952,mccluskey1956} finds all prime
implicants by iterative pairwise merging of minterms differing in one bit,
with don't-care cells included alongside on-set cells.  Essential prime
implicants cover on-set cells matched by no other implicant; greedy selection
completes the cover.  Two properties are machine-verified: every ON cell is
covered (\texttt{cc\_cover\_complete}) and no implicant matches any OFF cell
(\texttt{cc\_off\_avoiding}).

Flipped maps assign ON status to never-observed combinations, producing
forbidden-combination rules that encode negative-selection constraints.

\section{Information-Theoretic Measures and Direct Coupling Analysis}
\label{sec:theory-info-dca}

\subsection{Information-theoretic measures}
Shannon entropy per column quantifies evolutionary variability after excluding
gap characters.  Mutual information between columns $(i,j)$ measures
statistical dependence:
$\MI(i,j) = \sum_{a,b} P_{ij}(a,b) \log_2[P_{ij}(a,b)/(P_i(a)P_j(b))]$.
The Average Product Correction removes phylogenetic bias:
$\MIp = \MI - \bar{\MI}(i,\cdot)\bar{\MI}(\cdot,j)/\bar{\MI}$~\cite{dunn2008}.
Perplexity ratio $r(i,j) = \mathrm{PP}(j)/\overline{\mathrm{PP}(j|i)}$
captures near-determinism of the residue relationship.

\subsection{Mean-field DCA}
Direct Coupling Analysis models sequences through a Potts Hamiltonian whose
couplings are estimated by covariance inversion~\cite{weigt2009,morcos2011}.
Sequence reweighting at $\theta=0.2$ and pseudocount $\lambda=0.5$ regularise
the estimates.  The Frobenius norm of each coupling block provides a pair
score; APC correction removes phylogenetic bias.  Alignments deeper than
12{,}000 sequences are uniformly subsampled.

\section{Formal Verification in Lean 4}
\label{sec:theory-lean}

All load-bearing definitions are formalised in Lean~4~\cite{lean4} across two
packages totalling 118 theorems.  Key results include: Gray-code consecutive
adjacency, Hamming-distance symmetry, He-code pair counts at each distance,
separation-bin monotonicity, cell-index injectivity, implicant-cover
completeness and soundness, fully-fixed implicant uniqueness on the padded
grid, and padding safety.  Every theorem depends only on \texttt{propext},
\texttt{Quot.sound}, \texttt{Classical.choice}, and the
\texttt{native\_decide} evaluation axiom; none depends on \texttt{sorryAx}.

% ═══ Results: Cross-Protein Transfer ═══
\section{Universal Contact Circuits Transfer Across Proteins}
\label{sec:res-transfer}

We trained Quine--McCluskey-minimised Boolean circuits on consensus-residue
identities and separation distances from approximately 100 PSICOV150 proteins,
then evaluated on 50 entirely held-out proteins (maximum pairwise 5-mer
Jaccard index 0.029 across all $\binom{150}{2}$ pairs).  The identity-level
circuit achieved a mean top-$L/5$ precision of \textbf{0.179} ($5.1\times$
enrichment over base rate 0.035), a $1.9\times$ improvement over plain MI
(0.097), and a $15\times$ improvement over separation-only features (0.033).
Binary classification at circuit membership achieved MCC $\approx$ 0.095 with
enrichment of $2.5$--$2.8\times$.  Ranked precision was invariant across all
nine labelling-threshold combinations tested.

The minimised circuits produced interpretable rules.  Level-1 essential prime
implicants included hydrophobic$\times$aromatic packing at separations 6--21
and sulfur-mediated contacts at separations 6--11.  Level-2 vocabularies
included $\{M,F\}\times\{M,F,Y,W\}$ and $\{I,V\}\times\{I,V,F,W\}$,
corresponding to known aliphatic/aromatic core interactions rediscovered
without biophysical prior knowledge.

% ═══ Results: Gray Adjacency and Rank Fusion ═══
\section{Gray Adjacency Enrichment and Rank Fusion with DCA}
\label{sec:res-gray-fusion}

\subsection{Gray adjacency among native contacts}

Among 47{,}430 native-contact pairs across 150 PSICOV proteins, 22.83\% are
Gray-adjacent ($d_G = 1$), compared to a composition-controlled background of
19.37\%, yielding an enrichment of $1.178\times$ ($p = 1.1 \times 10^{-77}$,
two-sided exact binomial).  Bootstrap confidence intervals from protein-level
resampling span [22.27\%, 23.40\%].

A permutation test over random relabelings of the twenty amino-acid codes
demonstrated that this enrichment is specific to the He/Gray assignment:
random code permutations yield a mean enrichment of $0.961 \pm 0.095$, placing
the true encoding at $z = 2.51$ above the null distribution.  An earlier
analysis restricted to a single protein found the same point estimate but
failed to reach significance due to insufficient sample size ($n = 38$).

\subsection{Rank fusion improves DCA ranking}

The Karnaugh-map cell-rate prior improved global ranking (circuit AUC 0.658
vs.\ DCA AUC 0.583) while degrading top-$K$ sharpness.  We hypothesised that
rank aggregation would capture both properties.  Five fusion rules were
evaluated under identical protein-level cross-validation:

\begin{table}[ht]
	\centering
	\caption{Rank-fusion results over five-fold protein-level CV.}
	\label{tab:fusion-nat}
	\begin{tabular}{lcc}
		\toprule
		Method            & prec@L/5 (mean $\pm$ sd)   & AUC            \\
		\midrule
		rawDCA alone      & $0.169 \pm 0.125$          & 0.583          \\
		circH prior alone & $0.124 \pm 0.088$          & \textbf{0.658} \\
		\midrule
		\textbf{ranksum}  & $\mathbf{0.204 \pm 0.104}$ & 0.638          \\
		geometric mean    & $0.203 \pm 0.103$          & 0.634          \\
		product           & $0.199 \pm 0.108$          & 0.592          \\
		\bottomrule
	\end{tabular}
\end{table}

The rank-sum fusion achieved prec@L/5 = \textbf{0.204}, exceeding rawDCA
(Wilcoxon $p < 10^{-4}$) and circH ($p < 10^{-4}$).  Cohen's $d = 0.31$
with bootstrap 95\% CI $[0.017, 0.054]$.  Mechanistic decomposition showed
that all rescued pairs had above-median scores from both methods, confirming
coincidence-of-evidence amplification.  Adding plain MI as a third signal did
not improve further (three-way 0.201 vs.\ two-way 0.204, $p = 0.607$).

% ═══ Results: Complete Circuit (Positive + Flipped + Conservation) ═══
\section{The Complete Circuit Combines Positive, Flipped, and Conservation Channels}
\label{sec:res-complete}

We constructed a two-sided classifier over a shared feature space of
physicochemical groups, separation bins, and conservation-pair types
(10 bits, 1024 cells).  Two ternary tables were labelled from identical pooled
counts: the positive table marks cells whose contact rate significantly exceeds
the base rate, while the flipped table marks cells whose rate is significantly
below it.  A pair receives a definitive CONTACT call if its cell is ON in the
positive table and a definitive FORBIDDEN call if its cell is ON in the
flipped table; all other pairs are abstained.

On held-out proteins, the positive channel achieved precision 0.092
($2.6\times$ base) at recall 0.203.  The flipped channel achieved specificity
of \textbf{98.3\%}: forbidden-called cells showed exactly the predicted
$\leq$1.75\% contact rate on unseen proteins, confirming that negative-selection
rules generalise across protein families.  Combined definitive-call MCC was
$+0.169$ at 22\% coverage.

Pairs in which both alignment columns are fully conserved (entropy $< 0.001$)
showed a contact rate of 11.4\%, representing a $3.2\times$ enrichment.
However, a finer analysis using continuous minimum-pair entropy revealed that
absolutely invariant positions are depleted of contacts ($0.69\times$), while
near-conserved positions (entropy $0.001$--$0.05$) are enriched ($1.35\times$).
This non-monotonicity indicates that structural-core membership requires
residues to be constrained but not absolutely frozen; catalytic or functional
constraints can lock a position without placing it in a contact-rich
environment.

% ═══ Results: The Information Ceiling of a Feature Class ═══
\section{The Information Ceiling of a Feature Class}
\label{sec:res-ceiling}

Section~\ref{sec:res-transfer} reports that the universal circuit transfers to
held-out proteins at prec@$L/5 = 0.179$ against published DCA's $0.723$, and
Section~\ref{sec:res-reconstruction-negative} that this is insufficient for
template-free reconstruction.  Both are empirical statements about one method.
This section replaces them with a bound that holds for \emph{every} method
built on the same features, and thereby identifies where the shortfall lives.

The bound rests on an observation that is trivial and decisive: a circuit whose
inputs are a feature vector is a function of that feature vector, and therefore
cannot distinguish two candidate pairs that share it.

\begin{definition}[Cell predictor]
\label{def:cellpredictor}
Let $P$ be the candidate residue pairs of a protein, $y : P \to \{0,1\}$ the
native-contact labels, and $\varphi : P \to \mathcal{C}$ a feature map into a
finite set of \emph{cells}.  A \emph{cell predictor} is any
$\rho : \mathcal{C} \to \mathbb{R}$, inducing the score $\rho \circ \varphi$.
Pairs sharing a cell receive equal scores; their relative order is fixed by a
tie-break which carries no information about $y$ and is therefore modelled as
uniformly random.
\end{definition}

Every circuit in this paper is a cell predictor: a sum-of-products over
$\varphi$ is a $\{0,1\}$-valued function of the cell, and the smoothed
cell-rate scoring used in practice is a real-valued one.  The class of cell
predictors is strictly larger than the class of circuits, so a bound over it
binds the circuits a fortiori.

\begin{theorem}[Ceiling]
\label{thm:ceiling}
Write $n_c = |\varphi^{-1}(c)|$ and $r_c$ for the empirical contact rate in
cell $c$.  For any cell predictor $\rho$ and any $k$,
\[
\mathbb{E}\bigl[\operatorname{prec@}k(\rho \circ \varphi)\bigr]
\;\leq\; \frac{1}{k}\max\Bigl\{ \textstyle\sum_{c} t_c\, r_c \;:\;
0 \le t_c \le n_c,\ \sum_c t_c = k \Bigr\},
\]
the expectation taken over the tie-break, with the maximum attained by filling
cells in decreasing order of $r_c$.
\end{theorem}

\begin{proof}
Fix $\rho$.  Its induced ranking visits cells in decreasing $\rho$-order; write
$t_c$ for the number of pairs drawn from cell $c$ among the top $k$, so
$0 \le t_c \le n_c$ and $\sum_c t_c = k$.  Within a cell the tie-break is
uniform, so the $t_c$ selected pairs form a uniformly random $t_c$-subset of
$\varphi^{-1}(c)$ and the expected number of contacts drawn from $c$ is
$t_c r_c$ by linearity of expectation; hence
$\mathbb{E}[\operatorname{prec@}k] = k^{-1}\sum_c t_c r_c$.  Varying $\rho$
varies only the cell order, hence only the allocation $(t_c)$, which ranges
over prefix-filling allocations; relaxing to all allocations satisfying the
stated constraints can only increase the maximum.  That relaxation is a
fractional knapsack with unit weights, whose optimum sorts cells by $r_c$
descending---and this optimum is itself prefix-filling, so the relaxation is
tight.
\end{proof}

\begin{corollary}
\label{cor:nomethod}
No choice of minimisation algorithm, labelling threshold, don't-care policy,
cover-selection rule, or Boolean representation can raise the expected
precision of a predictor above the bound of Theorem~\ref{thm:ceiling} for its
feature map.  Improvement requires refining $\varphi$.
\end{corollary}

\begin{proof}
Each listed choice determines which cell predictor $\rho$ is obtained; none
changes $\varphi$.  The bound is uniform over $\rho$.
\end{proof}

Table~\ref{tab:ceiling} evaluates the bound on all 150 PSICOV targets at
$k = L/5$, with cells scored by their own pooled empirical contact rate---the
feature class handed the test labels.

\begin{table}[htbp]
\centering
\caption{Achievable prec@$L/5$ per feature map, 150 PSICOV targets.
Consensus residue codes; separation bins $[6,11)$, $[11,21)$, $[21,31)$,
$[31,\infty)$; mutual-information quartile computed per protein.}
\label{tab:ceiling}
\begin{tabular}{lrr}
\toprule
Feature map $\varphi$ & Cells & Achievable \\
\midrule
residue$_i$, residue$_j$, separation, MI quartile & 6390 & 0.2412 \\
residue$_i$, residue$_j$, separation              & 1600 & \textbf{0.1933} \\
physicochemical group$_i$, group$_j$, separation  & 256  & 0.1418 \\
residue$_i$, residue$_j$                          & 400  & 0.1383 \\
Gray Hamming distance, separation                 & 24   & 0.0789 \\
separation only                                   & 4    & 0.0598 \\
\midrule
\multicolumn{3}{l}{\emph{Measured on the same benchmark:}} \\
Trained circuit, held out (Sec.~\ref{sec:res-transfer}) & & 0.179 \\
Plain mutual information                                & & 0.097 \\
Published DCA~\cite{jones2012}                          & & 0.723 \\
\bottomrule
\end{tabular}
\end{table}

Three readings follow.  \textbf{The machinery is saturated}: the trained
circuit reaches $0.179$ against an achievable $0.1933$, or $93\%$ of its
feature class's ceiling, so by Corollary~\ref{cor:nomethod} the total remaining
headroom from every possible improvement in Boolean method is at most about
$8\%$ relative.  \textbf{The shortfall is informational}: the ceiling itself
sits $3.7\times$ below published DCA, and no Boolean construction on these
features can close that gap.  \textbf{Coarsening chemistry costs information}:
the eight-group physicochemical partition has a \emph{lower} ceiling than raw
residue identity, as it must, since merging cells can only reduce the bound.
Chemical classes buy interpretability, not accuracy.

This reframes Section~\ref{sec:res-reconstruction-negative}.  The failure to
meet the reconstruction criterion is not a deficiency of our minimiser or our
labelling; it is a property of consensus residue identity as an input.  Since
the bound depends only on the partition, the only lever is a finer partition
separating contacting from non-contacting pairs \emph{within} existing cells.
Adding a single crudely binned column statistic already lifts the ceiling from
$0.1933$ to $0.2412$, whereas richer residue vocabularies lift nothing---which
is the formal counterpart of the rank-fusion result of
Section~\ref{sec:res-gray-fusion}, where combining the circuit with continuous
DCA outperforms either alone.

% ═══ Results: Reconstruction Assessment and Negative Results ═══
\section{Reconstruction Assessment and Negative Results}
\label{sec:res-reconstruction-negative}

\subsection{Reconstruction is not achievable at current signal strength}
We assessed whether the circuits supply sufficient correct spatial constraints
for template-free structure determination by measuring recall (fraction of
native contacts recovered) and long-range precision.  Circuit recall@L = 0.028
versus published DCA's 0.218; median long-range precision@L/5 is 0.000 for
circuits versus 0.683 for PSICOV predictions.  Against the threshold of
median LR prec@L/5 $\geq 0.30$ and recall@L $\geq 0.10$ derived from
Jones et al.\ and the CONFOLD literature, the circuits do not meet the
criterion while published DCA does.

\subsection{Tension mechanism does not improve prediction}
An attention-style tension score, normalising each position's mutual
information by its strongest partner, did not improve circuit precision
(tension-augmented 0.154 vs.\ non-tension 0.165).  High-tension positions were
anti-correlated with structural hubs (median $\rho = -0.06$), indicating that
entropic attention concentrates on variable surface clusters rather than
buried cores.

\subsection{Iterative DI is unreliable on conserved-rich families}
Our mean-field implementation of iterative direct information produced
anti-predictive rankings (AUC $\approx 0.41$) on families with many fully
conserved columns, due to pseudo-inverse conditioning when variance approaches
zero for locked positions.  APC-corrected Frobenius norm was adopted as the
primary DCA score instead; this choice was validated against native contacts
and matches published mean-field benchmarks.

\subsection{Cross-dataset pilot}
A pilot evaluation on three EVmutation proteins produced inconclusive results
because correct parsing of the A2M alignment format requires handling
insert-state columns that do not map linearly to PDB residue numbers.
Structures were successfully fetched from RCSB and MI AUCs were positive for
two of three targets, but precision could not be meaningfully evaluated
without proper reference-frame mapping.  Full evaluation is deferred to
future work.

% ═══ Results: Boolean Logic over Conformational State ═══
\section{Boolean Logic over Conformational State}
\label{sec:res-contact-logic}

The results so far apply Boolean minimisation to \emph{sequence}: an alignment
column pair becomes a Karnaugh map, and its prime implicants describe which
residue combinations coevolution permits.  Section~\ref{sec:res-ceiling} shows
that this substrate is close to exhausted for contact prediction.  Here we
change the substrate rather than the method.

The motivation is that Boolean minimisation is productive elsewhere in biology
precisely where the object is a \emph{logical relation among variables observed
across many states}---gene-regulatory inference being the established case, in
which binarised expression time series yield logic rules via
Quine--McCluskey~\cite{kocevar2024sailor}.  A static residue string is not such
an object.  A structure-based model (SBM) trajectory~\cite{clementi2000sbm}
is: every native contact is either formed or broken in every frame, so the
trajectory is a binary state matrix and one may ask whether a contact's state
is a Boolean function of a few other contacts' states.

\subsection{Controlling the folding coordinate}
\label{sec:cl-control}

Every pair of native contacts is correlated for a trivial reason: when the
protein is folded most contacts are formed, when unfolded most are broken.  Any
naive coupling statistic measures this global mode and nothing else.  Three
devices control for it, all applied throughout.

First, $Q$, the fraction of native contacts formed, is quantile-binned into two
bits and supplied to the circuit as an \emph{explicit input}, so that a rule
reads ``at this folding level, contact $c$ is formed iff \dots'' rather than
``contact $c$ is formed when the protein is folded''.  Second, every rule is
scored against a \emph{$Q$-only baseline} built by the identical procedure from
the $Q$ bits alone; the reported quantity is always the gain over that
baseline, so a rule that merely restates the global mode scores zero.  Third, a
\emph{$Q$-stratified null} resamples contact columns within $Q$ strata,
preserving each contact's dependence on the folding coordinate and its marginal
frequency while destroying any specific contact-to-contact coupling.

Predictors are ranked by \emph{conditional} mutual information
$I(c; x \mid Q\text{-bin})$; plain mutual information would rank candidates by
how strongly each tracks the confounder.  Predictor selection is repeated
inside every cross-validation fold, so the \emph{choice} of predictors is held
out and not only the fitted rule.  Folds are contiguous frame blocks rather
than random splits.

\subsection{Coupling beyond the folding coordinate is present}

Across 480 proteins and 18{,}714 extracted rules, every circuit
runtime-verified sound and complete, the mean gain over the $Q$-only baseline
is $-0.0054$ while the $Q$-stratified null gives $-0.0248$.  The real
configuration exceeds its own null in \textbf{463 of 480 proteins}
(Wilcoxon $p = 4.3\times10^{-78}$).

The sign deserves comment.  On held-out frames the contact-based circuit is
marginally \emph{worse} than the two-bit $Q$ model, yet decisively better than
its own null.  Varying the predictor count identifies the cause
(Table~\ref{tab:cl-predcount}): the absolute gain degrades monotonically as the
truth table grows, while the real-minus-null difference stays flat at
$\approx +0.024$ and equally significant at every size.  That is a
variance-limited estimator, not absent signal---with 100 frames there are
$\approx 2.5$ frames per cell at three predictors.

\begin{table}[htbp]
\centering
\caption{Predictor count against estimator variance (60 proteins, 25 targets
each).  Absolute gain falls with table size; the real-minus-null difference
does not.}
\label{tab:cl-predcount}
\begin{tabular}{rrrrrl}
\toprule
Predictors & Cells & Gain vs.\ $Q$-only & Null gain & Real $-$ null & Real $>$ null \\
\midrule
1 & 8  & $+0.0033$ & $-0.0195$ & $+0.0228$ & 52/55 ($p = 2.8\times10^{-10}$) \\
2 & 16 & $-0.0002$ & $-0.0259$ & $+0.0257$ & 50/55 ($p = 6.0\times10^{-10}$) \\
3 & 32 & $-0.0023$ & $-0.0260$ & $+0.0237$ & 51/55 ($p = 3.5\times10^{-10}$) \\
\bottomrule
\end{tabular}
\end{table}

\subsection{The learned rules are spatially local}
\label{sec:cl-local}

The substantive finding is what the rules connect.  For each rule we compare
the contact-map $L_1$ distance $|i-k| + |j-l|$ between the target contact and
its selected predictors against predictors drawn at random from the
\emph{same protein's} varying contacts, so that protein size and contact
density cancel.

\begin{table}[htbp]
\centering
\caption{Selected predictors versus matched-random predictors, 480 proteins.
Adjacency is $L_1 \leq 2$ on the contact map.}
\label{tab:cl-locality}
\begin{tabular}{lrr}
\toprule
 & Observed & Matched random \\
\midrule
Fraction of predictors adjacent      & \textbf{0.2127} & 0.0170 \\
Enrichment                           & \multicolumn{2}{c}{\textbf{12.5$\times$}} \\
Proteins with observed $>$ random    & \multicolumn{2}{c}{\textbf{478 / 480}} \\
Wilcoxon                             & \multicolumn{2}{c}{$p = 3.3\times10^{-80}$} \\
Median contact-map distance          & \textbf{61.5} & 107.5 \\
Predictors sharing a residue         & \multicolumn{2}{c}{\textbf{27.2\%}} \\
\bottomrule
\end{tabular}
\end{table}

Conditional mutual information, with the global folding mode controlled out,
selects the \emph{neighbouring cells of the contact map}: adjacent rungs of the
same $\beta$-ladder, successive turns packed against the same partner.  More
than a quarter of selected predictors share a residue with their target.  This
is local cooperativity---the foldon picture~\cite{maity2005foldons}---recovered
from a discrete-logic direction without secondary structure being supplied.
Individual rules can be strong even where the population mean gain is not:

\begin{quote}\small
\texttt{1zjy}\quad contact $(1,28)$: accuracy $0.900$ vs.\ $Q$-only $0.480$, MCC $0.80$;\\
\phantom{\texttt{1zjy}}\quad predictors $(1,29), (2,28), (2,29)$, contact-map $L_1 = 1, 1, 2$.

\texttt{1kms}\quad contact $(2,129)$: accuracy $0.790$ vs.\ $0.380$, MCC $0.62$;\\
\phantom{\texttt{1kms}}\quad rule $c(2,129) = c(2,127) \wedge c(1,129)$, $L_1 = 2, 1$.
\end{quote}

\subsection{Longer trajectories confirm the variance diagnosis}
\label{sec:cl-dense}

The diagnosis in Table~\ref{tab:cl-predcount} makes a testable prediction:
more statistically independent configurations should flip the absolute gain
positive.  We re-simulated ten proteins for four times longer with twenty times
denser output ($2\times10^{6}$ steps saved every 250, against $5\times10^{5}$
every 5000) and repeated the identical analysis in three arms
(Table~\ref{tab:cl-dense}).

\begin{table}[htbp]
\centering
\caption{Three arms on the same ten proteins.  The middle arm decimates the
long run back to the \emph{original} 5000-step spacing, so its time resolution
is unchanged and only the number of independent configurations differs.}
\label{tab:cl-dense}
\begin{tabular}{lrrrrr}
\toprule
Condition & Frames & Spacing & Gain vs.\ $Q$-only & Null & Fraction $>0$ \\
\midrule
Original     & 100   & 5000 & $+0.0164$ & $-0.0549$ & 0.90 \\
Long-sparse  & 400   & 5000 & $\mathbf{+0.0628}$ & $-0.0299$ & \textbf{1.00} \\
Long-dense   & 8000  & 250  & $\mathbf{+0.0842}$ & $-0.0022$ & \textbf{1.00} \\
\bottomrule
\end{tabular}
\end{table}

Long-sparse is the informative arm: identical time resolution to the original,
four times the independent configurations, and the gain rises $3.8\times$ and
becomes positive in every protein ($+0.0464$, 10/10, $p = 0.00195$).  The
constraint was sampling, not method.  These ten proteins were selected as small
and strongly fluctuating, so their absolute levels are optimistic relative to
the 480-protein population; the paired comparisons are unaffected.

\subsection{Dense sampling licenses a kinetic reading}
\label{sec:cl-kinetic}

In the original trajectories the lag-1 autocorrelation of contact state
averages $0.009$ ($Q$: $0.130$): the 5000-step save interval exceeds the
correlation time, the frames are near-independent equilibrium samples, and no
$t \rightarrow t+1$ statement is licensed.  We measured this per protein rather
than assuming it, and made no kinetic claim on that data.

At 250-step spacing the contact lag-1 autocorrelation rises to $0.539$.  A
lagged analysis---predicting the target $k$ frames \emph{ahead} of the inputs,
with predictor selection lagged identically---therefore becomes meaningful
(Table~\ref{tab:cl-lag}).

\begin{table}[htbp]
\centering
\caption{Predicting a contact ahead of its inputs, ten densely sampled
proteins.}
\label{tab:cl-lag}
\begin{tabular}{rlrrl}
\toprule
Lag & Horizon & Gain vs.\ $Q$-only & Null & Real $>$ null \\
\midrule
0 & same frame  & $+0.0842$ & $-0.0022$ & 10/10 ($p = 1.95\times10^{-3}$) \\
1 & 250 steps   & $+0.0593$ & $-0.0023$ & 10/10 ($p = 1.95\times10^{-3}$) \\
2 & 500 steps   & $+0.0431$ & $-0.0024$ & 10/10 ($p = 1.95\times10^{-3}$) \\
4 & 1000 steps  & $+0.0288$ & $-0.0021$ & 10/10 ($p = 1.95\times10^{-3}$) \\
\bottomrule
\end{tabular}
\end{table}

Predictive power decays roughly exponentially, with a relaxation of
$\approx 900$ steps, but remains positive in every protein a thousand
integration steps ahead.  The state of three neighbouring contacts now carries
information about a target contact's future state beyond what the folding level
supplies---a statement about mechanism rather than only about the equilibrium
ensemble, and the first such statement this framework has been entitled to
make.

% ═══ Discussion ═══
\section{Discussion}
\label{sec:discussion}

We have demonstrated that the analytical chain from biological sequences
through Gray-code encoding, Karnaugh-map construction, Quine--McCluskey
minimisation, and prime-implicant extraction can be made machine-provable in
Lean~4, and that the resulting Boolean circuits carry genuine structural
information when applied to protein multiple sequence alignments.

The principal finding is that the Karnaugh-map cell-rate prior carries
orthogonal structural information demonstrable by rank-fusion with continuous
DCA scores.  The fused predictor achieves prec@L/5 = 0.204, exceeding both
the circuit alone (0.124) and DCA alone (0.169) at $p < 10^{-4}$.  This is
the first evidence that Boolean-minimisation features improve continuous
coevolutionary contact prediction.

The biological interpretation of learned circuit rules confirms known
structural principles.  Level-1 essential implicants encode hydrophobic-core
packing, aromatic stacking, and sulfur-mediated contacts in directly readable
sum-of-products form.  At Level-2 resolution, specific residue-pair
vocabularies correspond to aliphatic/aromatic packing interactions that
dominate the buried cores of globular proteins.

Several limitations constrain the scope of our conclusions.  Consensus-residue
features collapse within-column frequency variation into a single code,
discarding information that richer encodings might exploit.  Separation bins
use globally defined boundaries without per-protein length normalisation.
Cross-dataset generalisation was inconclusive on EVmutation due to A2M format
parsing complexity, leaving open the question of whether findings replicate
beyond PSICOV150.  The iterative direct-information computation from our
mean-field implementation remains unreliable on families with many fully
conserved columns; we adopted APC-corrected Frobenius norm as the primary DCA
score for this reason.

Three extensions follow naturally.  Replacing binary conservation flags with
continuous entropy bins as additional Karnaugh-map dimensions could capture
finer-grained structural information about the buried core.  Integrating
circuit-derived rules as input features to deep learning architectures would
combine the interpretability of Boolean logic with the representational power
of attention-based models.  Applying the complete pipeline to CASP
free-modelling targets would test whether the observed transfer generalises
to harder benchmarks with shallower alignments.

% ── Methods ──
% ═══ Methods: Dataset ═══
\section{Dataset}
\label{sec:methods-data}

The PSICOV150 benchmark~\cite{jones2012} provides 150 protein domains with
deep MSAs generated by three iterations of JackHMMER against UniRef100, together
with repaired native structures from the Protein Data Bank.  Alignment depths
range from 511 to 74{,}836 sequences; protein widths range from 50 to 266
residues.  The number of alignment columns equals the number of C-$\alpha$
atoms in each PDB file, verified empirically for 149 of 150 targets.

Native contacts use C$\beta$--C$\beta < 8$~\AA{} (C$\alpha$ for glycine) with
separation $\geq 6$, following CASP conventions~\cite{ezkurdia2009}.  Across
all targets, 47{,}430 native contacts exist among 1{,}675{,}886 candidate
pairs, yielding a base rate of 2.83\%.

Published PSICOV contact predictions from the supplementary data serve as
literature baseline.  Our measurement of these predictions yields prec@L/5 =
0.727 across all separations $\geq 6$, consistent with Jones et al.'s reported
value of approximately 0.73.

% ═══ Methods: Encoding and Feature Extraction ═══
\section{Encoding and Feature Extraction}
\label{sec:methods-encoding}

Each MSA is converted to a dense integer matrix using the He~2012 amino-acid
code (values 0--19), with gaps mapped to state~20.  Consensus residues are
computed as the majority non-gap code per column with lowest-code tie-breaking.
Candidate pairs include all $(i,j)$ with $j-i \geq 6$ and both columns having
valid consensus codes.

Pair features comprise: physicochemical group codes ($g_i, g_j \in \{0,\ldots,7\}$),
separation bins ($s \in \{0,1,2,3\}$ for $[6,11)$, $[11,21)$, $[21,31)$,
$[31,\infty)$), consensus residue codes, Gray-Hamming distance, mutual
information, and perplexity ratio.

Two circuit levels are constructed from pooled counts over training proteins:
Level-1 (256 cells) over group$_i \times$ group$_j \times$ separation-bin;
Level-2 (4$\times$1024 cells) over residue-identity pairs padded to
$32\times32$, one map per separation bin.  Cell labelling uses the ternary
rule: ON if contact rate $\geq T \times$ base rate and $n \geq n_{\min}$;
don't-care if $n < n_{\min}$; OFF otherwise, with defaults $T=2$,
$n_{\min}=30$.

% ═══ Methods: Circuit Training and Evaluation ═══
\section{Circuit Training and Evaluation}
\label{sec:methods-circuit}

Each labelled truth table is minimised using a popcount-grouped
Quine--McCluskey implementation with full don't-care support.  Every trained
circuit is asserted sound (no implicant matches any OFF cell) and complete
(every ON cell is covered) before evaluation.

Cross-protein evaluation uses protein-level cross-validation with folds
stratified by length quintile (seed~42).  All model parameters derive
exclusively from training proteins.  Ranked metrics include precision@L/5,
precision@L/2, precision@L, and AUC via Mann--Whitney $U$ with mid-rank tie
handling.

For the rank-fusion experiment, five rules are evaluated: rank sum,
$z$-score sum, product, geometric mean, and DCA-gated variant.  Paired
Wilcoxon signed-rank tests assess significance of improvements.

% ═══ Methods: Statistical Framework ═══
\section{Statistical Framework}
\label{sec:methods-eval}

Enrichment tests use exact one-sided binomial tests against
matched-separation random controls (2{,}000 draws per model, fixed seeds).
Spearman rank correlations assess monotone relationships.  Protein-level
bootstrap resampling (2{,}000 replicates) provides confidence intervals.
Wilcoxon signed-rank tests evaluate paired method comparisons without normality
assumptions.  Permutation tests over random code assignments quantify encoding
specificity.  Multiple-comparison awareness is maintained through median
reporting for skewed distributions and explicit Bonferroni-corrected thresholds
where noted.

All stochastic elements use fixed random seeds.  Re-running any stage produces
byte-identical output files, verified by MD5 checksums across multiple
independent invocations for each result category.

% ═══ Methods: Contact-State Logic ═══
\section{Methods: Boolean Logic over Contact States}
\label{sec:methods-contact-logic}

\subsection{Simulations and contact states}
Structure-based models~\cite{clementi2000sbm} were generated with
SMOG~2~\cite{noel2016smog2} as C$\alpha$-only topologies and integrated with
OpenSMOG~\cite{oliveira2021opensmog} on OpenMM~\cite{eastman2017openmm}, using a
Langevin middle integrator at reduced temperature $T = 0.99994$, time step
$0.0005$, collision rate $1.0$ and cutoff $1.1$~nm.  The production set
comprises 1{,}359 protein trajectories of $5\times10^{5}$ steps saved every
5000, of which the 500 lexicographically first with parseable topologies were
analysed and 480 yielded rules.

A native contact $(i,j)$ from the SMOG shadow map is \emph{formed} in a frame
when the C$\alpha$--C$\alpha$ distance is below $1.2$ times its distance in the
SBM reference structure; $Q$ is the fraction of native contacts formed, and
contacts whose formed-fraction lies outside $[0.05, 0.95]$ over the trajectory
are treated as constant and excluded as targets or predictors.

For the dense re-sampling of Section~\ref{sec:cl-dense}, ten proteins were
re-simulated with identical parameters for $2\times10^{6}$ steps saved every
250.  These 50--300 atom systems are small enough that OpenMM's CPU platform at
one thread ($\approx 12.7$k steps/s) outruns CUDA ($\approx 9$k), where
kernel-launch overhead dominates; the batch therefore ran one single-threaded
process per protein.  Six further candidates were dropped because their
topology directories lacked the OpenSMOG XML file.

\subsection{Rule construction}
Inputs to each circuit are the states of $n_{\mathrm{pred}} = 3$ selected
contacts followed by two bits of the quantile-binned $Q$, giving $2^5 = 32$
cells.  Predictors are the top-$n_{\mathrm{pred}}$ contacts by conditional
mutual information $I(c; x \mid Q\text{-bin})$, computed on the training frames
only and recomputed inside every fold.  Conditional rather than plain mutual
information is essential: plain MI ranks candidates by how strongly each tracks
the folding coordinate, which is the confounder being controlled.

Each cell is labelled by majority vote over its training frames---ON when the
target is formed in more than half, OFF otherwise---and cells holding fewer
than four frames become don't-cares, the standard treatment of unvisited states
in Boolean network inference and what allows the minimiser to generalise rather
than memorise.  The ternary table is minimised by exact Quine--McCluskey on the
sparse term set, and every returned cover is runtime-verified sound (no cube
covers an off-cell) and complete (every on-cell is covered) by cube
enumeration.

\subsection{Evaluation and controls}
Cross-validation uses five contiguous frame blocks rather than random splits:
where any temporal correlation survives the save interval, a random split would
place adjacent frames on both sides and inflate held-out accuracy.  Three
predictors are scored on the same held-out frames---the circuit, a $Q$-only
baseline built by the identical procedure from the $Q$ bits alone, and the
training-majority class---and the reported effect is always
$\mathrm{acc}_{\mathrm{rule}} - \mathrm{acc}_{Q\text{-only}}$.

The $Q$-stratified null resamples every contact column within $Q$ strata.  This
preserves exactly each contact's marginal frequency inside each $Q$ bin, hence
its entire dependence on the folding coordinate, while destroying which frames
pair which contacts.  A gain that survives this null is not explicable by the
global mode.

For the locality analysis, predictors are compared against contacts drawn at
random from the same protein's varying set, so that protein length and contact
density cancel; adjacency is $L_1 \leq 2$ on the contact map, where
$L_1 = |i-k| + |j-l|$ for contacts stored as sorted residue pairs.

Temporal scope is measured rather than assumed: the lag-1 autocorrelation of
$Q$ and of each varying contact is reported per protein, and the lagged
analysis of Section~\ref{sec:cl-kinetic} is applied only to the densely sampled
trajectories, where it exceeds $0.5$.

\subsection{Reproducibility}
Seeds are fixed throughout (base 20260905, offset per protein by index; folds
contiguous), and stage re-runs reproduce identical JSON.  The
Quine--McCluskey engine is shared with the sequence-circuit analysis and
carries 22 correctness gates, including agreement with an independent dense
implementation when no don't-cares are present and an independent audit of
implicant-hood and primality at 8, 10 and 12 variables.

We note one consequence of that audit for the wider pipeline.  The dense
reference minimiser retains only don't-cares at Hamming distance 1 from a real
on-minterm, a documented performance choice; where don't-cares cluster it
therefore under-merges, returning 187 cubes against 522 exact prime implicants
on a 10-variable test table, with every extra cube it reports strictly
contained in a true prime.  Rules obtained from it remain sound, but
prime-implicant counts and ``essential'' designations computed with it are
approximations rather than exact Quine--McCluskey quantities.

% ── References ──
\bibliographystyle{unsrtnat}
\bibliography{bib/references}

\end{document}
